\begin{auxlemma}[A single dual-standard translation of a Bessel period]
\label{lem:paired-lowering-factor}
Let \(\lambda_1>\cdots>\lambda_n\) be a marked ambient
Harish--Chandra parameter, \(u_1>\cdots>u_k\) the compact
parameter, and \(\mu=\lambda+\nu\) with
\(\|\nu\|_1+1\le q_0\). Take a nonzero Bessel model at \(\mu\)
and its normalized period maps as in
Proposition~\ref{prop:real-period-norm}. Let \(\kappa_s\) be the scalar
obtained by inserting the distinguished covector in the dual standard
tensor, projecting to the discrete series of parameter \(\mu-e_s\)
and applying its period map. In the same Haar measure at both parameters,
\[
 |\kappa_s|^2\frac{d_\mu}{d_{\mu-e_s}}
   =T_z R_s\le T_z B_0^4,\qquad
 R_s=\frac{\prod_{i=1}^k|\mu_s-\tfrac12-u_i|}
              {\prod_{t\ne s}|\mu_s-1-\mu_t|}.
\]
Here \(d_\mu\) is the formal degree and \(-T_z\) the squared
Hermitian norm of the distinguished negative vector.
\end{auxlemma}

\begin{proof}
\proofstep{Compute the projected period.}
The shifted summands in the dual Pieri rule have multiplicity one.
On the \(s\)-summand the split Casimir has eigenvalue
\(-\mu_s+(1-n)/2\), by subtracting the two factor Casimirs from
the tensor Casimir. Let \(P_s\) denote its spectral projection. In the
Capelli column determinant use the central polynomial \(Q(Y)\)
and its signed comatrix \(C(Y)\), with argument \(Y+(n-1)/2\).
The quantum comatrix identity gives
\[
 Q(Y)=\prod_t(Y+\mu_t),\qquad
 C(Y)P_t=\prod_{j\ne t}(Y+\mu_j)P_t .
\]
Indeed multiplication by the shifted split Casimir in the comatrix
identity multiplies the \(t\)-summand by \(Y+\mu_t\); polynomial
division gives the second equality, including at the roots.
In particular \(C(-\mu_s)=\prod_{t\ne s}(\mu_t-\mu_s)P_s\).
The determinant and minor conventions are those of
\cite[Proposition~2.7(4) and Remarks~2.8(1), 2.11]{mno-yangians}.

Use the ordered Witt basis consisting of the two null flags and
the middle block. In the Bessel functional the comatrix entry
between the outer frame vectors reduces to the compact \(k\)-row
minor. To see the factors, number the basis positions from
\(0\) to \(2r+k\). For \(r>0\) the cofactor deletes the last row
and the first column. Expand first in the first remaining row,
then in the last remaining column, successively down the two flags.
The nonzero root character entries outside the middle are
\[
 (a,a+1)\quad(a+1<r\text{ or }a>r+k),\qquad
 (r-1,r+k),\quad(r+k,r+k+1).
\]

The ordinary flag entries have modulus one and the two joining
entries each have modulus \(\sqrt{T_z}\). At each outer expansion
the displayed entry is the only one on which the functional survives.
In a column expansion the argument drops by one, from commuting
the last matrix unit to the left using
\([E_{ij},E_{ab}]=\delta_{ja}E_{ib}-\delta_{ib}E_{aj}\).
After the \(r\) column removals the compact argument at the root is
\[
 -\mu_s+\tfrac{n-1}{2}-r=-\mu_s+\tfrac k2,\qquad
 Q_{\rm compact}(-\mu_s+\tfrac k2)
       =\prod_{i=1}^k(-\mu_s+\tfrac12+u_i).
\]
The latter equality is the compact Capelli central character. For
\(r=0\) the cofactor is already the compact block, and insertion
and contraction of the negative vector give \(T_z\) directly.

The tensor Hermitian form restricts on the summand to its unitary
form times a sign \(\epsilon_s\). Thus the period of the reverse
contraction is \(\epsilon_s\overline{\kappa_s}\) times the shifted
period. Indeed pairing the two contractions on smooth vectors in
the period integral moves the adjoint across the Hermitian form,
and the inserted covector is fixed by the period subgroup.
On a vector with nonzero original period, the composite therefore
has modulus \(|\kappa_s|^2\). The cofactor evaluation gives
\[
 |\kappa_s|^2
 =T_z\frac{\prod_i|\mu_s-\tfrac12-u_i|}
              {\prod_{t\ne s}|\mu_s-\mu_t|},\qquad
 \frac{d_\mu}{d_{\mu-e_s}}
 =\prod_{t\ne s}\frac{|\mu_s-\mu_t|}{|\mu_s-1-\mu_t|}.
\]
The second identity follows by cancelling unchanged factors in
Harish--Chandra's Vandermonde formula at fixed Haar measure.
The admissible shifts remain regular, so all the denominators
are nonzero. Multiplication gives the asserted equality.

\proofstep{Pair the factors before estimating the quotient.}
In Definition~\ref{hyp:critical-level-weights} the upper entries
are \(v_1,\ldots,v_{k-1}\) and the remaining entries are
\(-b_1,\ldots,-b_m\). Use \(B_v=b_1\), \(q_0=km^5\) and the
spacing bound \(\gamma=nm^{10}\) from
Lemma~\ref{lem:marked-paired-moment-range}. Its selected row bounds give
\[
 \begin{gathered}
  0\le\lambda_i\le B_v/8\ (i<k),\qquad
  -2B_v\le\lambda_i\le-B_v\ (i\ge k),\\
  |u_i|,|\lambda_t|\le B_0,\qquad
  q_0\le\min(\gamma/8,B_v/4,B_0).
 \end{gathered}
\]
Strict interlacing and opposite parity at the compact and ambient
parameters give \(\lambda_i+\tfrac12\le u_i\) for \(i\le k\)
and \(u_i+\tfrac12\le\lambda_{i-1}\) for \(2\le i\le k\).
Put \(N_i=|\mu_s-\tfrac12-u_i|\),
\(D_t=|\mu_s-1-\mu_t|\) and \(\delta_t=|\nu_t|\).
Each \(N_i\le4B_0\), while \(D_t\) for \(t\ne s\) is a positive
integer and hence at least one.

If \(s<k\), assign a distinct denominator to each indicated numerator:
\[
 \begin{array}{c|c|c}
   i&t&A\rule{0pt}{2.5ex}\\ \hline
   2\le i<s&i-1&\lambda_t+1-\mu_s\\
   s+2\le i\le k&i&\mu_s-1-\lambda_t
 \end{array}
\]
In the first line \(u_i\le\lambda_{i-1}-1/2\) and
\(u_i\ge\lambda_i+1/2>\mu_s-1/2\).
In the second, \(u_i\ge\lambda_i+1/2\) and
\(u_i\le\lambda_{i-1}-1/2\le\mu_s-1/2\).
The separation of \(\lambda_s\) from \(\lambda_i\) in the first
line and from \(\lambda_{i-1}\) in the second is at least
\(\gamma\), so in both lines
\[
 N_i\le A,\qquad D_t\ge A-\delta_t,\qquad
 A\ge\gamma/2,\qquad 2\delta_t\le A .
\]

If \(s\ge k\), pair every \(2\le i<k\) with \(t=i-1\), and put
\(A=\lambda_t+1-\mu_s\). Here \(\lambda_t,\lambda_i\ge0\) and
\(\mu_s\le-B_v+q_0\), so \(N_i\le A\), \(D_t\ge A-\delta_t\),
\(A\ge B_v/2\) and \(2\delta_t\le A\).
For each paired factor, taking \(c=\gamma/2\) or \(B_v/2\)
according to the case gives
\[
 \frac{N_i}{D_t}\le\frac{A}{A-\delta_t}
      \le 1+\frac{2\delta_t}{A}
      \le e^{2\delta_t/c},\qquad
 \prod_{i\,\mathrm{paired}}\frac{N_i}{D_{t(i)}}\le e^{2q_0/c}\le e .
\]
The upper case omits at most \(1,s,s+1\), the lower case at most
\(1,k\). Unused denominators are at least one. Consequently
\[
 R_s\le
 \begin{cases}64e B_0^3&s<k,\\16e B_0^2&s\ge k.\end{cases}
 \qquad R_s\le B_0^4\quad(B_0\ge64e).
\]
The last threshold holds uniformly at the selected family scale.
\end{proof}
